\documentclass[../main]{subfiles}
\begin{document}
\chapter{Exposición bombas}
\img{images/12/bomba.png}{Bomba centrífuga.}

\info{
Contenido}{
\begin{enumerate}
    \item Pistón.
    \item Diafragma.
    \item Placas deslizantes o Paletas Deslizantes.
    \item Engranajes
    \item Centrífugas
    \item Axiales
\end{enumerate}
}

\section{Contenido de la presentación.}
\begin{enumerate}
    \item \emph{Principio de funcionamiento:} Explicar cada una de las partes que componen a la bomba, tanto la parte de la bomba como del motor que entrega la energía, con especificaciones técnicas y un ejemplo concreto. 
    \item Ciclo de bombeo: Explicar un ciclo de funcionamiento de la bomba, o el princípio físico que caracteriza su funcionamiento.
    \item \emph{Fotos}: Buscar al menos de 10 fotos de este tipo de bombas.
    \item \emph{Historia:} Historia de su fabricación, primeras aplicaciones en la industria, inventor creador, y primeros uso de este tipo de bomba.
    \item \emph{Características Técnicas}: consumo de energía potencia, altura,altura de aspiración, presión, caudal de funcionamiento, presiones negativas, positivas.
    \item \emph{Mantenimiento}: Explicar como se realiza el mantenimiento de esta bomba.
    \item \emph{Ventajas y desventajas:} Analizar las ventajas y desventajas del uso de este tipo de bomba. 
    \item Mercado actual:
    \begin{itemize}
        \item ¿Dónde se compra?
        \item ¿Cuál es el precio?
        \item ¿Que características tiene cada diferentes  modelo comercial?
    \end{itemize}
    \item \emph{Conclusiones:} Resumen con conclusiones generales después de la presentación de la clase.
    \item \emph{Preguntas:} Se responderán las preguntas que los alumnos tengan con respecto a la presentación.
    \item Subir archivo pdf con el informe y powerpoint (PPT), a classroom con el nombre de los integrantes.
\end{enumerate}


\end{document}
